\chapter{Discussion}
\label{sec:discussion}

%%%%%%%%%%%%%%%%%%%%%%%%%%%%%%%%%%%%%%%%%%%%%%%%%%%%%%%%%%%%%%%%%%%%%%%%%%%%%%%%
\section{Comparison with Related Work}
\label{sec:discussion:relatedwork}
% Compares as precisely as the literature allows, and states where a direct numerical
% comparison is impossible because no prior work reduces AIGs for this task.

\subsection{Positioning Against Optimizability Prediction Work}
\label{sec:discussion:relatedwork:prediction}
% Protocol must be matched before numbers are compared: most published figures sit in
% leakier protocol cells (\ref{sec:relwork:ml4eda:protocols}), and the RQ1a inflation
% factor (\ref{sec:results:rq1:protocol}) is what translates between them. SynthNet's
% collapse on this split is consistent with that audit rather than an implementation
% artifact.

\subsection{Positioning Against Graph Reduction Work}
\label{sec:discussion:relatedwork:reduction}
% Retention at matched compression against what coarsening-for-GNN studies report,
% and whether H1's contraction advantage appears in their non-DAG results too.

\subsection{Where Direct Comparison Is Not Possible}
\label{sec:discussion:relatedwork:incomparable}
% No prior work reduces AIGs for this task; states which nearest comparisons remain
% and what each one does and does not control for.


%%%%%%%%%%%%%%%%%%%%%%%%%%%%%%%%%%%%%%%%%%%%%%%%%%%%%%%%%%%%%%%%%%%%%%%%%%%%%%%%
\section{Interpretation of Findings}
\label{sec:discussion:interpretation}
% Explains the results rather than restating them, grouped by theme so that RQ2 and RQ3
% can be discussed as the single trade-off they describe.

\subsection{Feasibility of Optimizability Prediction}
\label{sec:discussion:interpretation:feasibility}
% What \ref{sec:intro:rq1} implies about whether structure alone carries the signal.

\subsection{What Reduction Buys and What It Costs}
\label{sec:discussion:interpretation:tradeoff}
% \ref{sec:intro:rq2} and \ref{sec:intro:rq3} together. Which point on the Pareto
% front is actually worth taking, and is the offline reduction cost recovered?

\subsection{Why Domain-Informed Adaptation Did or Did Not Help}
\label{sec:discussion:interpretation:domain}
% \ref{sec:intro:rq4}. If the AIG-aware heuristics do not win, what that says about
% where the predictive signal lives, distinguishing ``domain knowledge does not help
% here'' from ``these particular heuristics were weak''.

\subsection{Structural Generalization and Deployment}
\label{sec:discussion:interpretation:generalization}
% \ref{sec:intro:rq5}. Whether train-cheap/infer-full is a viable deployment story.

\subsection{Unexpected Results}
\label{sec:discussion:interpretation:unexpected}
% Anything that contradicted the hypotheses, including negative results. Better
% surfaced deliberately here than left for a reader to notice.


%%%%%%%%%%%%%%%%%%%%%%%%%%%%%%%%%%%%%%%%%%%%%%%%%%%%%%%%%%%%%%%%%%%%%%%%%%%%%%%%
\section{Practical Implications}
\label{sec:discussion:implications}
% Concrete recommendation: given a memory budget and an accuracy tolerance, which
% reduction should a practitioner reach for? This is the payoff of the Pareto front.


%%%%%%%%%%%%%%%%%%%%%%%%%%%%%%%%%%%%%%%%%%%%%%%%%%%%%%%%%%%%%%%%%%%%%%%%%%%%%%%%
\section{Limitations}
\label{sec:discussion:limitations}
% Limitations discovered or imposed, distinct from the scoping choices declared in
% \ref{sec:intro:rqs:scope}.

\subsection{Reproducibility}
\label{sec:discussion:limitations:reproducibility}
Every stochastic component of training is seeded, so an individual run repeats exactly. That
is a weaker guarantee than it appears. Each configuration was trained once, so no run-to-run
variance is available and every accuracy difference reported in Chapter~\ref{sec:results}
is a point estimate. This bears hardest on RQ4, where the expected effect is small: a gap of
a few percentage points between a domain-informed method and its control cannot be
distinguished from seed noise on the evidence collected here. The efficiency comparisons are
better placed, being paired per graph with confidence intervals and a signed-rank test, but
accuracy has no equivalent treatment.

Three further dependencies limit exact reproduction. Timing and memory measurements are
specific to the cluster hardware and to the enabled compilation and mixed-precision
settings, so they should be read as relative comparisons rather than as portable absolutes.
The labels themselves depend on the version of the synthesis tool used to generate them; a
different version would produce a different corpus, and the version was not pinned. The
corpus is also not uniformly regenerable. One of the four generation scripts is stochastic,
and its seed was drawn per graph and never recorded
(\ref{sec:apx:algorithms}), so the quarter of the tier-1 set it produced could be
reconstructed only in distribution, not exactly. Those graphs are training inputs rather
than label sources, so the targets this thesis reports are unaffected, but the corpus itself
is reproducible only in part.

\subsection{Scalability}
\label{sec:discussion:limitations:scalability}
This is the limitation closest to the motivation and it should be stated without hedging.
The argument for this work rests on industrial circuits of millions to billions of nodes.
The evaluation runs on a corpus bounded at roughly a third of a million nodes per graph,
because RQ1 requires a full-graph baseline and a baseline that does not fit in memory cannot
be measured. The design that makes the study interpretable is the same design that keeps it
away from the scale it argues about.

What does extrapolate: the relative ordering of reduction methods by offline cost, and the
qualitative shape of the memory and throughput trade-offs, since these follow from
properties of the algorithms and of message passing rather than from the corpus size. What
does not: the absolute compression figures, since several methods' compression is a property
of circuit structure rather than a parameter, and larger industrial circuits may have
different redundancy profiles; and any claim that a method which retains accuracy here would
retain it at a scale where the model sees a proportionally smaller fraction of each graph.

\subsection{Generalizability}
\label{sec:discussion:limitations:generalizability}
Three axes are held fixed, each narrowing what the conclusions cover.

\emph{One synthesis algorithm.} Results characterise optimizability under
\texttt{Orchestrate}. Graphs optimized by the other three scripts exist and are unused. If
the reduction rankings hold across all four, the finding is about AIG structure; if they do
not, it is about \texttt{Orchestrate}, and nothing here can tell the two apart.

\emph{One encoder.} A four-layer edge-aware GCN+ was used throughout. Depth is not
incidental to these results: methods were coupled to it by design, with refinement depth set
equal to the number of layers. A deeper model would merge less under colour refinement and
would have a different receptive-field profile, so the ranking could shift.

\emph{One target.} Node reduction is the only quantity modelled. Delay and power are the
objectives synthesis actually trades against, and nothing here speaks to them.

\subsection{Reliability and Validity}
\label{sec:discussion:limitations:validity}
\emph{Construct validity.} Node reduction is a proxy for area, and area is one of three
objectives. A circuit whose node count barely moves may still have been substantially
improved in delay. The label measures what the tool removed, not how much better the result
is, and ``optimizability'' should be read in that narrower sense throughout.

\emph{Internal validity.} Splitting by design was chosen to prevent a base graph and its own
optimized descendants from straddling the train/test boundary. Whether it fully succeeds is
testable rather than assumed, and the test is run as RQ1a
(\ref{sec:results:rq1:protocol}): the gap between a design-level split and a random split
measures directly how much of the score is design recognition rather than structural
inference.
% TODO: once the RQ1a numbers land, replace the hedge above with the measured inflation
% factor, and state whether per-design degradation was uniform or concentrated. If
% concentrated, a second design fold is the mitigation, at one training run.

\emph{Construct validity of the comparison itself.} Hyperparameters were tuned once on the
full-graph configuration and held fixed. This protects the comparison from confounding
retention with tuning effort, but it means reduced configurations run at settings chosen for
a different input distribution. That biases against them, and therefore bounds how
strongly any negative result about reduction can be stated.

\subsection{Limitations of the Reduction Methods Tested}
\label{sec:discussion:limitations:methods}
The reduction methods are not a representative sample of their literatures, and the
summarization family in particular was scoped down deliberately: it runs one domain-specific
method, one general GNN-aware method and one exact method, against a random control. The
classical, guarantee-bearing spectral branch of the coarsening literature is not among them,
so this thesis reports no trained result for it.

The domain-informed adaptations are simple. Edge-span weighting, level slicing, gate-role
filtering and reconvergence-bounded merging are direct encodings of structural intuitions,
not sophisticated methods. If they fail to beat their generic counterparts, that bounds
these heuristics, not the idea of domain-aware reduction. That is a
distinction~\ref{sec:results:rq4:adequacy} carries, and it should not be quietly dropped
when the conclusion is written.

Finally, the boundary between structural compression and functional optimization is argued
rather than measured. The claim that merging structurally equivalent nodes does not leak the
label while merging functionally equivalent nodes would is well founded, but the negative
control that would demonstrate it empirically, running FRAIG-style functional reduction
\citep{mishchenko2005fraig} as a summarizer and observing the label leak, was not run. The
super-node content ablation is in the same position:
member statistics are computed but not yet consumed by the encoder, so the finding in the
literature that super-node content matters is cited rather than tested here.


%%%%%%%%%%%%%%%%%%%%%%%%%%%%%%%%%%%%%%%%%%%%%%%%%%%%%%%%%%%%%%%%%%%%%%%%%%%%%%%%
\section{Ethical Considerations}
\label{sec:discussion:ethics}
\paragraph{Compute and energy cost.}
The work has a real environmental footprint: generating a corpus of several million AIGs
required sustained CPU cluster time, and the training sweeps required substantial GPU time
on top. The tension is worth stating rather than eliding: a study whose purpose is to
reduce the compute cost of training spent a large amount of compute establishing how. The
defence is that the cost is incurred once and the artifacts are stored and reusable, so a
result showing which reductions preserve the signal saves compute for every subsequent user;
but that is a claim about future savings against a certain present cost.
% TODO: quote the actual GPU-hours and node-hours consumed rather than describing them
% qualitatively. The figures make this paragraph honest instead of gestural.

\paragraph{Benchmark provenance and licensing.}
The corpus is derived from published benchmark circuits, and their licences govern what may
be redistributed. Every non-synthetic netlist was received through the OpenABC-D
distribution, so the copies used are governed by its BSD 3-Clause terms, with the EPFL
suite's MIT license applying in addition (\ref{sec:method:data:sources}). A dataset release
must therefore carry both attribution notices. The ISCAS-85 and MCNC circuits attach no
formal terms; their inclusion rests on the same decades-long redistribution practice that
OpenABC-D itself relies on.

\paragraph{Accessibility, and who benefits.}
Lowering the memory required to train circuit models cuts in two directions. It widens
access, since research that currently requires large multi-device clusters could run on modest
hardware, which matters for groups without industrial compute budgets. It also lowers the
cost of capability for those who already have that budget. The honest position is that this
work makes an existing capability cheaper rather than creating a new one, and that the
widening effect is the more likely of the two to matter at this scale.

\paragraph{Dual use.}
Logic synthesis tooling is general-purpose infrastructure. Improvements to it apply to
whatever circuits are being designed, and this work does nothing specific to any application
domain. No mitigations are proposed because none are specific enough to be meaningful; the
consideration is recorded because omitting it would be a worse answer than naming it.


%%%%%%%%%%%%%%%%%%%%%%%%%%%%%%%%%%%%%%%%%%%%%%%%%%%%%%%%%%%%%%%%%%%%%%%%%%%%%%%%
\section{Outlook}
\label{sec:discussion:outlook}
% A few sentences pointing forward; the full treatment is
% \ref{sec:conclusion:futurework}.
